\section{Additionner des nombres relatifs}
\begin{propbox}
    \bb{Règle 1} : Pour \bb{additionner} deux termes de \hl{même signe}:
    \begin{itemize}
        \item On ajoute les \bb{distances à zéro}.
        \item On garde le même signe.
    \end{itemize}
\end{propbox}


\begin{exemplebox}
    \begin{itemize}
        \item $A = 5 + 7 = \answer{12}$ 
        \item $B = -5 + (-10) = \answer{15}$ 
    \end{itemize}
\end{exemplebox}

\begin{propbox}
    \bb{Règle 2} : Pour \bb{additionner} deux termes de \hl{signes contraires}: 
    \begin{itemize}
        \item On \bb{garde le signe du plus grand}. 
        \item On soustrait leur \bb{distances à zéro}. (\it{"Le plus grand moins le
            plus petit"})
    \end{itemize}
\end{propbox}


\begin{exemplebox}
    \begin{itemize}
        \item $A = 5 + (-7) = \answer{-2}$ 
        \item $B = -5 + 10 = \answer{+5}$ 
    \end{itemize}
\end{exemplebox}

\begin{methodebox}
\begin{enumerate}
    \item Je regarde les signes.
    \item Je choisis la bonne règle.
    \item Je calcule les distances à zéro.
    \item Je donne le bon signe.
\end{enumerate}
\end{methodebox}
